
\section{Conclusion}

In this work, we propose extreme token compression, named \textbf{\name}, to address the long video understanding problem by alleviating the large number of tokens. At the token level, our \name framework targets on compressing each frame to \emph{one token}. Observing the failures of heuristic compression adopted by previous methods, we propose the key insight of letting the LLM layers \emph{learn} the \emph{progressive} compression of video tokens. This compression strategy, called LP-Comp, effectively improves the token efficiency and enables denser frame sampling. Our compression at the frame level further enhances the information contained in the tokens by selecting the frames relevant to the user queries. Based on the internal attention scores of LLM layers, we further propose to split the whole video into shorter segments to mitigate the position bias issues. The resulting \emph{question-conditioned} compression, QC-Comp, curates a selective set of frames for the LLM to process. Collectively, our proposed strategies enable \name to explore the extreme of video token compression via data-efficient fine-tuning and significantly improve the base VideoChat-Flash on multiple long video understanding benchmarks.

\noindent \textbf{Limitations and Future Work.} Under a limited budget, we primarily explore the extreme token compression with VideoChat-Flash-2B and LLaVA-Next-Video via data-efficient fine-tuning. Therefore, potential future work includes adapting our \name framework to larger VLMs and potentially integrating the learnable compression into the pre-training stage of the VLM. In addition to the scale and diversity of experiments, another future work lies in training with the question-conditioned frame selection, which could be jointly optimized with the LLM layers to achieve even larger compression ratios than presented in the paper. 

\noindent \textbf{Broader Impacts.} This work presents advancements in vision-language modeling, with a focus on improving long video understanding. While our contributions are primarily foundational, we acknowledge potential societal risks associated with vision-language models, including misuse of disinformation, privacy concerns, and the amplification of biases present in training data. We encourage future work to include fairness assessments and responsible release strategies. Where applicable, safeguards should be considered to mitigate unintended harms.